\section{関係式・potential・選んだ周期構造}\label{sec:period-stack}
$V_i=\kk^4$ ($i\in\mathbb Z/4$) を頂点$i$から$i+1$への矢空間とし、
$G=\prod_i\mathrm{GL}(V_i)$、$\overline G=\prod_i\mathrm{PGL}(V_i)$とおく。
標準AS条件とは左右Koszul解消の型$(1,4,6,4,1)$、Gorenstein parameter 4、
行Hilbert級数$(1-t)^{-4}$を意味する。族では各成分の局所自由性を要求する。

\begin{proposition}[任意の周期を保持する関係式模型]\label{prop:relation-period}
$\prod_{i=0}^3\operatorname{Gr}(6,V_i\otimes V_{i+1})$内の標準AS部分関手を
$\mathcal R^{\rm AS}$とする。商stack$[\mathcal R^{\rm AS}/G]$は、
4周期同型$\tau:C\xrightarrow{\sim}C(4)$を\emph{指定した}標準AS正則
$\mathbb Z$-代数のstackと同値である。
\end{proposition}
\begin{proof}
$C_{i,i+1}$ ($0\le i<4$)の基底を選び、$\tau$によって全添字へ延長する。
二次関係の核が$R_{i+4}=R_i$を満たすので四つのGrassmannianの点となる。
逆に四つの関係を整数列上へ反復し、自由経路代数をその二次イデアルで割れば$C$を得る。
$\tau$を保つ同型は四つの基底変更で尽くされる。相対的には矢空間をfpqc局所で自明化して
同じ構成を行い、ベクトル束・積・関係の降下を用いる。AS部分への制限も両方向で一致する。
これはAS locus全体が有限型代数的stackであるという主張ではない。
\end{proof}

\paragraph{符号と微分の規約.}
$\mathcal W=V_0\otimes V_1\otimes V_2\otimes V_3$とし、最初の矢を最後へ送る回転を$c$とする。
$w\in\mathcal W$から、異なる始点成分を持つ
$W(w)=w-cw+c^2w-c^3w$を作る。これは$cW=-W$を満たす四次superpotentialである。
代数は長さ2の経路微分で$D(W,2)=\kk Q/(\partial_pW:|p|=2)$と定義する。
BSW\cite[Theorems 6.2, 6.8]{BSW}はKoszul (twisted) CY代数と、対応する
解消複体が完全な (twisted) superpotentialの対応を与える。
通常の符号なしcyclic quotientへ$W$を送ると零になるので、その意味のpotentialと混同しない。

\begin{proposition}[任意の4周期には非捩れpotentialは対応しない]\label{prop:period-counterexample}
命題\ref{prop:relation-period}の対象全体を、非捩れ四次superpotentialのモジュライと
無条件に同一視することはできない。
\end{proposition}
\begin{proof}
$C=\check{\kk[x_0,x_1,x_2,x_3]}$とし、$\sigma=\operatorname{diag}(2,1,1,1)$から
$\tau=\check\sigma\circ(\text{添字の4移動})$を選ぶ。
最初の四つの矢空間を通常の$V$と同定し、その後の基底を$\tau$で運ぶと
\[
 R_0=R_1=R_2=\wedge^2V,\qquad R_3=(1\otimes\sigma^{-1})\wedge^2V.
\]
この最後の空間は$\wedge^2V$でない。例えば$x_0\otimes x_1-\frac12x_1\otimes x_0$
が属する。
もし同じ四窓が非捩れpotentialを持てば、始点0の最高syzygyテンソルは
\[
 (\wedge^2V\otimes V^{\otimes2})\cap
 (V\otimes\wedge^2V\otimes V)\cap
 (V^{\otimes2}\otimes\wedge^2V)=\wedge^4V
\]
に属する。従って非零スカラーを除いて交代テンソルであり、その巡回微分が与える
$R_3$も$\wedge^2V$となって矛盾する。標準AS性自体は元の多項式$\mathbb Z$-代数から保証される。
\end{proof}

正しい比較は、非捩れpotentialに対してはFrobeniusに適合する周期を用いること、
任意の周期を扱う場合はNakayamaの捩れとtraceの変更も含むtwisted potentialを用いることである。
幾何学的点で適合周期の存在を示すだけでは、任意の\emph{選択}を分類したことにならない。
また$[U/\mathrm{GL}_4^4]$と$[U/\mathrm{PGL}_4^4]$は異なる。
射影potentialでは$(\mathbb G_m)^4$が全点を固定し、前者にはその慣性が残る。
本稿で7次元というときは後者、または対応する一般有理商を指す。

\section{R(1,3)の四頂点表示と有限型の一般開部分}\label{sec:r13-explicit}
四つの矢空間を一時的に$V=\langle t,x,y,z\rangle$へ同定する。
完全対称化を$\operatorname{sym}$、四変数の完全交代テンソルを$\operatorname{Alt}(txyz)$とする。
\[
 F=f_3(x,y,z)+tq_2(x,y,z)+t^2\ell_1(x,y,z)
 \quad\bmod\kk t^3
\]
は固定した中心直線に対して19係数を持つ。次のテンソルを定める。
\begin{equation}\label{eq:r13-potential}
 w_F=\operatorname{Alt}(txyz)
       -3(1-c+c^2-c^3)\bigl(t\operatorname{sym}(F)\bigr).
\end{equation}
式中の回転は、いったん全slotを$V$に同定して行う。得られたテンソルの各slotを
$V_0,V_1,V_2,V_3$へ戻して$\mathcal W$の元とみなす。
その四頂点superpotentialは上の$W(w_F)$である。

\begin{proposition}\label{prop:potential-presentation}
式\eqref{eq:r13-potential}の長さ2の微分が生成する関係空間は各始点で
\begin{equation}\label{eq:r13-relations}
 [t,x]=[t,y]=[t,z]=0,\qquad
 [y,z]=\operatorname{sym}F_x,\quad
 [z,x]=\operatorname{sym}F_y,\quad
 [x,y]=\operatorname{sym}F_z
\end{equation}
の6次元空間に等しい。$t^3$を$F$へ加えても$w_F$は変わらない。
$f_3$の平面三次が滑らかな範囲では、この族はnoetherian、Koszul、標準Hilbert型の
4次元CY代数を与える。
\end{proposition}
\begin{proof}
全二階微分を$V^{\otimes2}$の語で展開する。
$(\partial_{xy}-\partial_{yx})w_F=2[t,z]$とその巡回式から三つの中心関係を得る。
$\partial_{tx}w_F$などは、これらの中心関係で割った後に
$[y,z]-\operatorname{sym}F_x$などとなる。
逆包含は、各語$x_ix_j$を可換二次式$x_ix_j+\Pi_{ij}/2$へ写すことで直接確認できる。
ここで$\Pi_{xy}=F_z$など、$\Pi_{ti}=0$である。
この写像の核は式\eqref{eq:r13-relations}そのものである。
各式は$F$の19係数についての多項式恒等式であり、付属コードは独立な19文字を残して展開する。
$F=t^3$に対する補正は$(1-c+c^2-c^3)t^4=0$である。

AS条件について、$t$を係数として三変数potentialの低次摂動とみなす。
最高次代数はMori--Smith\cite[Theorem 4.2(1)]{MS}の$A_{f_3,1}$であり、
滑らかな$f_3$では3次元CYである。低次cyclic potentialのPBW定理
\cite[Proposition 7.3.1(ii)]{EG}により、active次数の随伴は$A_{f_3,1}\otimes\kk[t]$。
各$t$の特殊化のPBW基底を持ち上げると$\kk[t]$上自由であり、$t$は中心正則一次元となる。
中心一次拡大のKoszul AS上昇\cite{LBS}から標準4次元AS正則性を得る。
式\eqref{eq:r13-potential}と完全Koszul解消を使い、BSWによりCY性を得る。
パラメータ基底上では各固定次数の表示行列のrankが一定で、被約な基底上の局所自由性と
解消の各内部次数の完全性が従う。必要なら有限個の基底chartへ制限する。
\end{proof}

\subsection{七係数chartと全四頂点の接空間}
\begin{equation}\label{eq:seven-parameter}
 F_{\boldsymbol a}=a(x^3+y^3+z^3)+bxyz
       +t(x^2+c_2y^2+c_3z^2)
       +t^2(d_1x+d_2y+d_3z)
\end{equation}
を用いる。パラメータは$(a,b,c_2,c_3,d_1,d_2,d_3)$の7個である。
可逆なHesse変換とshearによる一般形としても理解できるが、以下では厳密な横断行列で
支配性を示すので、一意な大域正規形を仮定しない。

$Z_6\subset\PP(\mathcal W)=\PP^{255}$を二つのcyclic flattening
$M_{01|23},M_{12|30}$ ($16\times16$) がともにrank6となる局所閉schemeとする。
rank条件は全7次minorの消滅と適切な6次minorの非消滅で課す。
対応する二次関係は、これらおよび転置flatteningの像である。

\begin{theorem}[四頂点成分の一般開部分]\label{thm:r13-open}
$\overline G=(\mathrm{PGL}_4)^4$とする。
$\overline G\cdot[w_F]$の像の既約閉包$\mathcal Z_{13}$は$Z_6$の67次元既約成分であり、
一般点で正則である。ある非空開集合$N\subset\mathbb A^7$と
$\overline G$-不変な非空開部分$O_{13}\subset\mathcal Z_{13}$が存在して、
\[
 \overline G\times N\longrightarrow O_{13},\qquad(g,\boldsymbol a)\longmapsto g[w_{F_{\boldsymbol a}}]
\]
はエタール全射となる。$O_{13}$の全代数は標準Koszul AS正則である。
従って$[O_{13}/\overline G]$は7次元の滑らかな有限型DM stackで、$N$がエタールatlasである。
\end{theorem}
\begin{proof}
式\eqref{eq:seven-parameter}の点
\[
 \boldsymbol a_*=(3,5,2,3,1,2,4)
\]
を取る。最高三次のHesse判別条件は$9^3+5^3=854\ne0$であり、前命題が適用できる。
この点のflatteningを$M$とするとrankは6。
rank条件の一次式は、左核$L$と右核$K$の各組に対する
$L^{\mathsf T}\dot M K=0$である。両flatteningについて合計200本、未知数256個の行列を$T$とする。
四つの$\mathfrak{gl}_4$による軌道微分を$O$、7係数変分を加えた行列を$S_7$とする。
付属の有理数上の行列証明書により
\begin{gather*}
 \operatorname{rank}T=188,\qquad
 \operatorname{rank}O=61,\qquad
 \operatorname{rank}S_7=68,\qquad TS_7=0.
\end{gather*}
従ってprojective接空間は$256-188-1=67$次元であり、族の像の次元も少なくとも$68-1=67$。
源$\overline G\times\mathbb A^7$の次元が$60+7=67$なので、局所次元と接空間次元が一致する。
従って$Z_6$はこの点で正則であり、像の閉包はその既約成分である。
射影軌道は$61-1=60$次元なので、stabilizerのLie代数は零。
両方が滑らかな67次元空間となる近傍へ制限すると、全射な接写像は同型となり、射はエタールとなる。
そのエタール locus の像は開である。源の$N$を縮小し、その$\overline G$飽和を$O_{13}$とする。
源の全点のAS性は前命題で保証されている。相対的な次数成分と解消はエタール忠実平坦降下で
$O_{13}$へ降りる。stabilizerのLie代数零条件も開に制限できるため商はDMである。
\end{proof}

これは単にPymの一頂点7次元族の係数を数えた議論ではない。
\emph{全256テンソル方向}で接空間を計算し、異なる4周期表示の方向も含めて成分性を認定している。
ただし閉包の全境界におけるAS条件、正規性、正規化、有理不変量の生成元は未決定である。
具体的なエタール群oidは
\[
 \mathcal G_7=\{(\boldsymbol a,\boldsymbol a',g):
        [w_{F_{\boldsymbol a'}}]=g[w_{F_{\boldsymbol a}}]\}
       \rightrightarrows N
\]
であり、行列のliftでは$(G_0\otimes\cdots\otimes G_3)w_{F_{\boldsymbol a}}
=u w_{F_{\boldsymbol a'}}$、$u\ne0$、$\det G_i\ne0$という有限式になる。
一般次数や有限stabilizerの位数をまだ算出していないので、$N$を粗モジュライそのものとは呼ばない。
GL商のstack次元は$67-64=3$で、4次元スカラー慣性をrigidifyすると7となる。


\subsection{中心一次拡大の次数付き同型とアフィン商}
\label{sec:r13-affine-quotient}
以下では四頂点の任意周期の同型ではなく、\emph{一頂点の次数付き代数}
$B_f$の同型を考える。これは大域Zhang同値の主張とは異なる。
$U=\kk^3$、$\mathcal F_{19}=\{f\in\kk[U]_{\le3}:f(0)=0\}\simeq\mathbb A^{19}$
とし、$f_3$が滑らかな平面三次を定める非空開集合を
$\mathcal F_{19}^{\rm sm}$とする。
$f=f_3+f_2+f_1$を$t$で斉次化した$F_f$に対して
式\eqref{eq:r13-relations}で定まる次数付き代数を$B_f$と書く。

\begin{theorem}[R(1,3)中心拡大の次数付き同型]\label{thm:r13-affine}
$f,f'\in\mathcal F_{19}^{\rm sm}$とする。
$B_f$の次数一次の中心は$\kk t$であり、
$B_f$と$B_{f'}$が次数付き代数として同型である必要十分条件は、
ある$(A,b)\in\operatorname{GL}_3\ltimes U=\operatorname{Aff}_3$
によって
\begin{equation}\label{eq:affine-action}
 f'(u)=\rho_{A,b}(f)(u):=
 \frac{f(Au+b)-f(b)}{\det A}
\end{equation}
となることである（同型の向きを反転すれば群作用の逆を用いる）。
任意の次数付き同型は、全次数一次生成元の共通スカラー倍を除けば
$t\mapsto t$、$u\mapsto Au+b t$の形に一意に表される。
従って次数の共通$\mathbb G_m$-スカラーをrigidifyした
\emph{次数付き表示族}の同型群oidは
\[
 [\mathcal F_{19}^{\rm sm}/\operatorname{Aff}_3]
\]
で表される。
定理\ref{thm:r13-open}の有理数証明点を含むある非空不変開部分では
stabilizerが有限であり、この商は滑らかな7次元Deligne--Mumford stackである。
\end{theorem}
\begin{proof}
$B_f$の二次関係は、$\wedge^2\kk^4$から$\operatorname{Sym}^2\kk^4$
へのgraphをなすので、二次部分は完全対称化により
$\operatorname{Sym}^2\kk^4$と線型同型である。
一次元$\ell=a_1x+a_2y+a_3z+ct$が中心なら、
$\ell$と$x,y,z$の交換子はこの対称化表示で
$a\times\nabla f_3$を最高active次数に持つ。
$a\ne0$ならこの三つが零であることは、
適当な線型座標で$f_3$が一変数の三乗になることを強制し、
平面三次の滑らかさに矛盾する。
従って$Z(B_f)_1=\kk t$である。
同じ計算は族でも成り立つ。一次中心を定める有限自由加群間の
交換子写像は、$\kk^3=\langle x,y,z\rangle$上で各幾何学的ファイバーに
階数3を持つ。よって基底上で局所的に三次小行列式を可逆にでき、
一次中心は$\OO_Tt$という直線束であり、任意の底変更と可換する。
従って以下の同型の記述はfpqc局所的な族についても適用できる。

次数付き同型$\psi:B_f\to B_{f'}$は$\psi(t)=c t'$を満たす。
全次数一次元を同じ$c^{-1}$倍する次数付き自己同型を
ターゲットに合成して$\psi(t)=t'$に規格化できる。
このとき$\psi(u)=Au'+b t'$で、$A$は可逆である。
$B_{f'}$の交換子を使えば、列ベクトル表示で
\[
 \begin{pmatrix}[\psi(y),\psi(z)]\\
                 [\psi(z),\psi(x)]\\
                 [\psi(x),\psi(y)]\end{pmatrix}
 =(\det A)A^{-T}\,\operatorname{sym}(\nabla F_{f'}).
\]
一方、$B_f$の関係を$\psi$で送ると、右辺は
$\operatorname{sym}(\nabla F_f(Au'+bt',t'))$である。
完全対称化は線型変換に両立するので、
$\nabla_{u'}F_f(Au'+bt',t')=(\det A)\nabla_{u'}F_{f'}(u',t')$
が必要十分である。
従って両斉次三次式は$t'^3$の倍数を除いて等しく、
$t'=1$で定数項を零へ規格化すると式\eqref{eq:affine-action}となる。
逆に同式のもとでは六つの二次関係が全て保たれるので、
$u\mapsto Au'+bt'$は次数付き同型を与える。
これにより同型の全てと共通スカラーによるrigidificationを得る。

$\dim\mathcal F_{19}^{\rm sm}=19$、
$\dim\operatorname{Aff}_3=12$である。
一般stabilizerが有限であることは次の簡単な点で直接検査できる。
\[
 f_*=x^3+y^3+z^3+x^2+2y^2+3z^2+x+2y+4z.
\]
$\operatorname{Aff}_3$のLie代数元$(M,\beta)$による
$f$の一次変化は、定数項を除いて
\[
 Df(u)(Mu+\beta)-(\operatorname{tr}M)f(u)
\]
である。
$f_*$の三次部分でこれが零なら、
$x_i^2x_j$の係数から$M$の非対角成分は零、
$x_i^3$の係数から$M=s\,\mathrm{id}$となる。
二次部分の係数から
$\beta=(s/3,2s/3,s)$であり、一次部分の$x$係数は
$-2s+2s/3=-4s/3$となる。
従って$s=0$、$M=\beta=0$である。
この点のstabilizerは標数零の有限型代数群でLie代数が零なので有限。
有限stabilizer条件の非空不変開へ縮小し、
滑らかな商の次元$19-12=7$を得る。
\end{proof}
\begin{remark}
Pym\cite[\S3.4]{Pym}および
Le Bruyn--Smith--Van den Bergh\cite[\S3]{LBS}の
中心拡大の一般表示と整合する定理であり、
\emph{graded algebra}の同型類を正確に記述したものである。
四頂点の任意周期付きZ代数の全同型を
$\operatorname{Aff}_3$で尽くしたとは主張しない。
四頂点potentialの商とgraded表示商を一般開上で同型にするには、
余分なZhang同一視を排除する別の比較が必要である。
\end{remark}

\begin{proposition}[八次臨界値多項式のアフィン共変性]
\label{prop:critical-discriminant-covariance}
$f\in\mathcal F_{19}^{\rm sm}$のJacobian環
$J_f=\kk[u_1,u_2,u_3]/(\partial_1f,\partial_2f,\partial_3f)$と
八次多項式$\Delta_f(\lambda)=\det(\lambda+m_f\mid J_f)$
を考える。
$f'=\rho_{A,b}(f)$とすると、アフィン変換$u\mapsto Au+b$は
$\operatorname{Spec}J_{f'}$と$\operatorname{Spec}J_f$を同型にし、
\begin{equation}\label{eq:critical-discriminant-covariance}
 \boxed{\quad
 \Delta_{f'}(\lambda)
 =(\det A)^{-8}\,
   \Delta_f\bigl((\det A)\lambda-f(b)\bigr)
 \quad}
\end{equation}
が成立する。これは非被約な臨界点schemeも含めて成り立つ。
特に、一般の8個のnodalな特異メンバーの対応は、
次数付き中心拡大$B_f$の同型に対して自然である。
\end{proposition}
\begin{proof}
式\eqref{eq:affine-action}を微分すると
$\nabla f'(u)=(\det A)^{-1}A^T\nabla f(Au+b)$となる。
従って二つのJacobianイデアルはアフィン座標変換で対応する。
対応する有限長8の環上では、$f'$の乗法は
$(m_f-f(b)\,\mathrm{id})/\det A$に移る。
その特性多項式を取れば式\eqref{eq:critical-discriminant-covariance}を得る。
この議論は固有値の相異性を使わず有限長環上で行うので
非被約な臨界点も含む。
一般の8つの通常二重点については
定理\ref{thm:pencil-critical}を適用する。
\end{proof}

\subsection{任意の周期を持つ関係式stackとの一般点での比較}
\begin{proposition}\label{prop:period-formal-rigidity}
$\check B$のFrobenius標準4周期点を取り、二次関係$R\subset B_1^{\otimes2}$の
二因子射影stabilizerが有限であるとする。
この点の全Artin標準AS四窓変形は、射影非捩れpotentialで一意に表示される。
従ってこの一般点では、任意周期付き関係式stackに潜在的な追加変形方向はない。
商に使うGLまたはPGLおよびスカラーrigidificationは両側で揃える。
\end{proposition}
\begin{proof}
Artin底$T$上の四窓のKoszul双対$E_T$は、長さ$0,\dots,4$で
$(1,4,6,4,1)$型の相対Frobenius代数である。
このことは標準AS条件の線型解消、または固定内部長の有限自由bar複体の底変更から従う。
最高次数の四つのtraceを選び、特殊ファイバーでは通常の意味で対称なtraceに規格化する。
そのNakayama自己同型$\nu_T$は各頂点を固定し、特殊ファイバーで恒等となる。
逆双対の次数一写像は隣接する二次関係を保存する。
小拡大の核$I$において$\nu=1+I D$と書くと、その射影像は特殊ファイバーの二因子
stabilizerのLie代数に入る。これは零なので、帰納的に各矢空間上の$\nu_T$はスカラー$c_i$である。
Nakayamaは最高次数traceを固定するから$\prod_{i=0}^3c_i=1$。
四つのtraceを単元$t_i$倍すると$c_i$は$(t_i/t_{i+1})c_i$へ変わり、積が1なので全て1にできる。
さらに頂点$i$のtraceを$(-1)^i$倍すればgraded対称traceとなる。
これを四次テンソルに双対化し、完全対合を用いれば関係空間はその二階微分となる。
二つのこのtraceは隣接完全対合により共通スカラーだけ異なるので、射影potentialは一意である。
この構成は関係式の同型と両立する。
\end{proof}
二因子stabilizerの有限性は第\ref{sec:surface-witness}節に記した$60\times32$行列のrank30によって、同じ$F_*$で直接確認される。
従って命題\ref{prop:period-counterexample}の大域的な不一致と、
$R(1,3)$一般点での形式的な一致は両立する。

\section{反標準環と認識部分のambient射}
次数一生成の次数付き環$A$に対し、その\emph{二次包絡}を
\[
 Q(A):=T_\kk(A_1)\Big/
 \big\langle\ker(A_1\otimes_\kk A_1\longrightarrow A_2)\big\rangle
\]
と定める。これは二次関係のデータを抽出する構成であり、
AS正則性は定義だけからは従わない。

\begin{proposition}[反標準環とambient候補]\label{prop:intrinsic-ring}
定義\ref{def:weak-cubic}に対し
\[
 \Gamma_{\rm ac}=\bigoplus_{m\ge0}\Hom(\OO,H^m\OO)
\]
は積も含めて内在的である。通常のArtin--Zhang再構成の有限性仮定の下で
$(\mathcal A,\OO,H)$は$(\qgr\Gamma_{\rm ac},\pi\Gamma_{\rm ac},(1))$から回復される。
$\Gamma_{\rm ac}$が次数一生成なら$Q(\Gamma_{\rm ac})$は一意な二次ambient候補である。
\end{proposition}
\begin{proof}
積は$f\in\Hom(\OO,H^m\OO)$、$g\in\Hom(\OO,H^n\OO)$に対して
$H^m(g)\circ f$である。Serre同値は特徴付けにより同値で保存され、
その同型の選択の変更はorbit algebraの同型を与える。
AZ amplenessと有限性の下でsection algebra再構成\cite{AZ}を適用する。
二次包絡は長さ二の合成核だけから定義される。
\end{proof}

内在的公理（rank-one torsion-free性とexceptionalityを含む）だけから、次数一生成や二次包絡のAS正則性が従うかは未証明であり、本稿では保留する。
本稿ではその一般的な反例を与えてはいないが、これらの含意は未証明である。
そこで内在的stackの\emph{超曲面認識部分関手}$\mathfrak C^{\rm hyp}$を、
$A=\Gamma_{\rm ac}$が次数一生成、$B=Q(A)$が標準AS4、
$I_3=\ker(B_3\to A_3)$がrank1でその生成元が左右正則、という条件で切る。
Hilbert級数と長さ4の核の次元4から左右乗法の像が一致し、
正規性と全高次核の生成が従って$A=B/(G)$となる。
族の局所自由性と分裂核から任意の底変更とも両立する。
従ってこの部分には実際の射
\begin{equation}\label{eq:ambient-map}
 \Phi:\mathfrak C^{\rm hyp}\longrightarrow\mathfrak P^{\rm AS,pt}_3,
 \qquad(\mathcal A,\OO)\longmapsto Q(\Gamma_{\rm ac})
\end{equation}
がある。ここでは源の自然同型と対応する頂点スカラーを処理したpointed targetを用いる。
rawな$\mathrm{GL}_4^4$商を用いる比較では、源にも同じ周期表示・スカラーのデータを加えた
presentation stackを用い、その正規incidenceの射を取る。内在的な源とrawなtargetを
そのまま比較して同じstack次元・同じファイバー次元を主張しない。
次数付きorbit algebraからの選択周期は次数shiftの4乗である。
非捩れpotentialのtargetを使う場合は適合周期へ規格化する。
この二つの選択は一般には同じ対象でなく、$R(1,3)$のCY表示上では適合する。

\section{AOUからの有限なambient構成と一般点の支配性}\label{sec:aou-ambient}
AOUの関係式族は、三つのblockを持つacyclic quiverの代数
$\Lambda$で、9頂点、18矢、9本の長さ2の関係を持つ。
長さ3の経路がなく、関係間のoverlapがないので、通常の相対quiver解消は長さ2で完全である。
$\Lambda$のrankは$9+18+18=45$である。
従って相対inverse Serre bimodule
$\Theta_\Lambda=\operatorname{RHom}_{\Lambda^e}(\Lambda,\Lambda^e)$を有限複体で計算できる。
$P_0$を頂点$(0,0)$の右射影対象とし、
\begin{equation}\label{eq:finite-serre}
 K_m=\operatorname{RHom}_\Lambda
       (P_0,P_0\mathbin{\otimes^L_\Lambda}\Theta_\Lambda^{\otimes^L m}[2m])
       \quad(0\le m\le4)
\end{equation}
と置く。これは係数環上の有限perfect複体である。

\begin{proposition}[低次数からの代数的候補射]\label{prop:aou-candidate}
滑らかな可換標識付き三次曲面を含み、相対inverse Serre複体の底変更が成立するAOUのある開部分上で、
$K_m$は次数0に集中し、そのrankは順に$1,4,10,19,31$となる。
$V=H^0K_1$、$R=\ker(V^{\otimes2}\to H^0K_2)$により、
有限Grassmannianの商への代数的射
\[
 \Psi_2:M_{\rm AOU}^{\rm low}\longrightarrow
 [\operatorname{Gr}(6,16)/\mathrm{PGL}_4]
\]
が得られる。これはまだ全域でAS性を結論する射ではない。
\end{proposition}
\begin{proof}
可換点では$P_0=\OO_X$、$\mathsf S^{-1}[2]=(-)\otimes\OO_X(1)$なので、
$K_m=R\Gamma(X,\OO_X(m))$である。$0\le m\le4$の高次コホモロジーは消える。
有限個のperfect複体に開性と底変更を適用し、rankと次数0集中を保つ開部分を取る。
$V^{\otimes2}\to H^0K_2$の全射性もこの点で成り立つので開に制限でき、核はrank6。
積はSerre orbitの合成から定まり、基底変更に両立する。従って所望の商への射がある。
\end{proof}


\begin{proposition}[一般幾何学的点でのinverse Serre orbit環]\label{prop:aou-orbit-generic}
AOU\cite[\S9, Theorem 9.4]{AOU}の一般位置条件
（滑らかな楕円曲線、無限位数のtranslation、十分一般の六点）を満たす
楕円nonupleに付随する九頂点有限代数を$\Lambda$、
頂点$(0,0)$の例外対象を$P_0$とし、
$H=\mathsf S^{-1}[2]$とする。積を保つ次数付き環同型
\[
 \bigoplus_{m\ge0}\Hom_{D^b(\Lambda)}(P_0,H^mP_0)
 \ \simeq\ F
\]
が存在する。ここで$F=P/(G)$はVan den Bergh\cite[Chapter 12]{VdB}の
反標準環であり、$P$は中心一次元$t$を持つ四次元AS正則Koszul代数、
$G$は正則正規三次元である。
この主張は幾何学的点での導来圏内の同型であって、
アーベルheartの一致や相対底変更を含まない。
\end{proposition}
\begin{proof}
AOU (9.44)--(9.45)の例外列の先頭は構造対象$\OO$であり、
(9.55)--(9.57)の導来同値は対応する$P_0$を
$\pi F\in\qgr F$へ送る。同論文Lemma 9.2により
$D^b\qgr F$のSerre関手は$(-1)[2]$なので、
$H$は次数shift$(1)$に対応する。
したがって左辺は
$\bigoplus_{m\ge0}\Hom_{\qgr F}(\pi F,\pi F(m))$
と積を保って同型になる。
$P$が四次元AS正則で$G$が正則三次元であることから、
$0\to P(-3)\xrightarrow{\cdot G}P\to F\to0$によって
$F$のgraded depthは少なくとも3である。
従って$H^0_{\mathfrak m}(F)=H^1_{\mathfrak m}(F)=0$であり、
Artin--Zhangのsection環の完全列\cite{AZ}を用いれば
$\Hom_{\qgr F}(\pi F,\pi F(m))\simeq F_m$が$m\ge0$で成立する。
\end{proof}

\paragraph{generic pointに用いる既知の入力.}
AOU\cite[Theorem 1.1, \S9, Theorem 9.4]{AOU}と
Okawa--Ueda\cite[Theorem 6.4]{OU2025}は、一般の標識付きデータを楕円曲線と六点blow-upで記述する。
命題\ref{prop:aou-orbit-generic}によりAOUの導来同値とSerre関手から、
一般幾何学的点でのinverse Serre orbit環は$F$に一致する。
同論文 (9.6)--(9.10) とVan den Bergh\cite[Chapter 12]{VdB}によれば
$F=P/(G)$で、$P$は中心一次拡大AS4、$G$は正則正規三次元。
幾何学的点での同型は命題\ref{prop:aou-orbit-generic}で証明したが、
相対底変更・形式変形への比較は残っている。
全標識付き成分やコンパクト化全体には拡張しない。

一般のSklyanin頂部では三次中心が一次元である。
正規化自己同型を$Gx=\sigma(x)G$と書くと、中心正則一次元$t$から$\sigma(t)=t$。
$P/(t)$上で$\overline G$は中心正則なので$\overline\sigma=1$となる。
よって$\sigma(x_i)=x_i+c_i t$である。
Pym/LBSの一般Jacobian座標で二次関係を比較すると
$\sum_i c_i\partial_i\partial_j f_3=0$であり、滑らかな$f_3$から$c_i=0$。
従ってこの一般部では$G$は中心的である。
標準の中心持上げと低次数中心の消滅により$Z(P)_3=\kk\Omega\oplus\kk t^3$となる。
一般の中心拡大の型がPymの$R(1,3)$であることは\cite[\S3.4]{Pym},\cite{LBS}の入力である。

\begin{theorem}[一般AOU点からのambientの像]\label{thm:aou-dominance}
命題\ref{prop:aou-orbit-generic}の一般位置に制限する。
さらに(i)この幾何学的点での復元が相対inverse Serre orbit環の底変更と両立し、
(ii)可換標識付き点からの形式変形の8方向をAOUの関係式表示がすべて実現し、
低次数写像の微分比較が適用できる、と仮定する。
この比較が成立する一般開部分では命題\ref{prop:aou-candidate}は
AS4 ambientへの有理射を与え、$R(1,3)$の7次元成分を支配する。
対応する中心三次対への有理像は8次元であり、その一般ファイバーは有限である。
従ってambientへの有理射の一般ファイバーは1次元である。
その有限次数を$|W(E_6)|$と同定したわけではない。
\end{theorem}
\begin{proof}
命題\ref{prop:aou-orbit-generic}により一般幾何学的点のorbit環は$P/(G)$であるから、二次包絡は$P$に一致する。
従って候補射の一般値は$R(1,3)$に属し、三次核もrank1となる。
これらは有限な式とrank条件なので、低次数の有理射として定義される。
支配性を以下で示した後、定理\ref{thm:r13-open}のatlasの逆像へ制限する。

像の次元を単なるパラメータ差から推測せず、truncationの微分を使う。
可換点においてAOUの8方向が全Abel変形を与えることは、定理の仮定(ii)の内容である。
反標準環の次数$0,1,2,3$の変形が、基底変更を除いて自明なら、その全次数変形も自明である。
実際、Artin底で次数一生成性は各次数のNakayamaにより持ち上がる。
自由テンソル代数からの核$J$は各次数で平坦であり、
$J/(VJ+JV)$の次数$>3$は特殊化で零なのでNakayamaにより零である。
従って核は次数2と3で生成され、低次数の自明化は環全体を自明化する。
AZ再構成によってAbel変形も自明となる。
よって低次数射の微分は単射であり、曲面同型類の像の次元は8である。
$A_1$の射影基底を加えると源の次元は$8+15=23$で、低次数像も23次元となる。

一般点ではこの低次数像は中心三次対であり、固定ambientに対する中心直線は$\PP^1$。
従ってambient像の次元は少なくとも$23-1=22$。
一方、中心拡大型ambientの座標付き像は22次元であるから等号。
一般の射影自己同型が有限である範囲で$\mathrm{PGL}_4$を割り、次元7と1次元一般ファイバーを得る。
四頂点potential側への移行は定理\ref{thm:r13-open}および
命題\ref{prop:period-formal-rigidity}の一般開部分を用いる。
\end{proof}

\begin{remark}[有理射と弱い対象の全域の射]
式\eqref{eq:finite-serre}は有限な行列・tensor積・核の計算による明示的レシピであり、
8個のトーラス座標から7個の正規形座標への展開済み有理式はまだ計算していない。
一般点の$\qgr F$はAOUの完全強例外列との導来同値によりsmooth properであり、
Serre公式とAZの定理により定義\ref{def:weak-cubic}を満たす。
従って式\eqref{eq:ambient-map}の像はこの成分の幾何学的一般点を含み、
その閉包の余次元は0である。
これは認識部分における\emph{一般点の支配性}であり、AOUの全開集合がそのまま
相対smooth properな族を与えることや、全境界を含む正則射、全内在的成分での認識定理は含まない。
\end{remark}
